\addcontentsline{toc}{chapter}{\textbf{Anexo 1}} 
\lhead{Anexo 1}
\rhead{\newtitle}
\cfoot{\thepage}
\renewcommand{\headrulewidth}{1pt}
\renewcommand{\footrulewidth}{1pt}
\chapter*{Anexo 1}


\noindent \textbf{Modelo estimación de GHI en condiciones de cielo claro utilizado}.


En este anexo se presentan las ecuaciones que definen el modelo \textit{ARGP-v2}, una versión 
mejorada del modelo \textit{ARG-P} propuesto originalmente por Salazar et al.~(2010) para la 
estimación de la irradiancia global horizontal en condiciones de cielo claro (GHI\textsubscript{cc}) 
en sitios de gran altitud.  
El modelo ARGP-v2 utiliza datos del servicio McClear de CAMS para recalibrar el índice 
representativo de claridad en función de la altura sobre el nivel del mar.

\section*{1. Ecuación general del modelo}

La irradiancia global horizontal en cielo claro estimada mediante el modelo se define como:

\begin{equation}
\mathrm{GHI_{cc}} = G_0 \, K_{trP} \, M_a^{-1}
\label{eq:argp_general}
\end{equation}

donde:
\begin{itemize}
  \item $G_0$: irradiancia extraterrestre horizontal.
  \item $M_a$: masa de aire óptica relativa corregida por presión.
  \item $K_{trP}$: índice representativo de claridad corregido por presión.
\end{itemize}

\section*{2. Masa de aire óptica relativa corregida por presión}

La masa de aire óptica utilizada se basa en la formulación de Kasten corregida por presión:

\begin{equation}
M_a = \left( \frac{P}{1013} \right)
\left[
\cos(z) + 0.50572 \left( 96.07995^\circ - z \right)^{-1.6364}
\right]^{-1}
\label{eq:ma_kasten}
\end{equation}

donde $z$ es el ángulo cenital solar y $P$ la presión atmosférica local en hPa.

\section*{3. Modelo ARG-P original}

El índice representativo de claridad corregido por presión ($K_{trP}$) en función de la altura $A$ 
está dado por:

\subsection*{Sitios con altura menor a 1000 m}

\begin{equation}
K_{trP} = 0.7570 + 1.0112 \times 10^{-5}\, A^{1.11067}
\label{eq:ktrp_low}
\end{equation}

\subsection*{Sitios con altura mayor o igual a 1000 m}

\begin{equation}
K_{trP} = 0.7000 + 1.6391 \times 10^{-3}\, A^{0.55000}
\label{eq:ktrp_high}
\end{equation}

\section*{4. Cálculo del RMSE para calibración}

Para ajustar el modelo ARGP-v2 se buscó minimizar el error cuadrático medio definido por:

\begin{equation}
\mathrm{RMSE} = 
\sqrt{
\frac{1}{M}
\sum_{i=1}^{M}
\left(
\mathrm{GHI_{cc}^{ARGP\_v2}}(i)
-
\mathrm{GHI_{cc}^{McClear}}(i)
\right)^2
}
\label{eq:rmse}
\end{equation}

donde $M$ es la cantidad de datos considerados.

\section*{5. Modelo ARGP-v2}

El modelo ARGP-v2 mantiene la estructura del ARG-P, pero reemplaza el índice de claridad por una 
nueva función ajustada en función de la altura:

\begin{equation}
\mathrm{GHI_{cc}^{ARGP\_v2}} = 
G_0 \, K_{trP\_v2}(A)\, M_a^{-1}
\label{eq:argpv2_general}
\end{equation}

La expresión ajustada del nuevo índice de claridad es:

\begin{equation}
K_{trP\_v2}(A) = 0.649 + 0.020\, A^{0.28}
\label{eq:ktrpv2}
\end{equation}

\section*{6. Ecuación final del modelo ARGP-v2}

Sustituyendo \eqref{eq:ma_kasten} y \eqref{eq:ktrpv2} en \eqref{eq:argpv2_general}:

\begin{equation}
\boxed{
\mathrm{GHI_{cc}^{ARGP\_v2}} =
G_0 \left[ 0.649 + 0.020\, A^{0.28} \right]
\left(
\frac{P}{1013}
\left[
\cos(z) + 0.50572 \left( 96.07995^\circ - z \right)^{-1.6364}
\right]^{-1}
\right)^{-1}
}
\label{eq:argpv2_final}
\end{equation}

Esta expresión constituye la formulación final del modelo ARGP-v2 para estimar la irradiancia 
global horizontal en condiciones de cielo claro.

